\section{结论}
\label{sec:conclusion}

本文面向标注稀缺和未知采集域条件下的小麦器官语义分割，提出统一框架 \METHOD。该方法不再将伪标签筛选、域泛化和高分辨率推理视为相互独立的步骤，而是以细长器官的拓扑可靠性贯穿训练与推理。\TRPL 将茎秆伪标签解耦为稳定中心线、可靠区域和不确定边界，在保留有效细结构监督的同时抑制噪声；\TCPM 从拓扑核心构建类别--域原型，并通过留一域模拟学习未知域中的稳定器官表示；\BAZR 根据熵、骨架端点和内部尺度差异选择少量困难窗口进行高分辨率修正。

在 GWFSS 少标注跨域协议和完整数据地域协议上的实验从总体 mIoU、茎秆 IoU、拓扑完整性和推理效率四个方面验证了所提框架。完整模型预计达到 \placeholder{验证 mIoU} 和 \placeholder{未知域测试 mIoU}，并将密集多尺度滑窗的约 \placeholder{前向次数} 次分割器前向计算降低至 $1+\placeholder{K}$ 次。后续工作将研究根据图像复杂度动态确定精细化预算，并将拓扑可靠监督扩展到其他具有遮挡和断裂特征的作物器官分割任务。
